\documentclass{article} % For LaTeX2e
\usepackage{iclr2026_conference,times}

% Optional math commands from https://github.com/goodfeli/dlbook_notation.
\input{math_commands.tex}

\usepackage{hyperref}
\usepackage{url}
\usepackage{amsmath,amssymb}
\usepackage{booktabs}
\usepackage{tabularx}
\usepackage{array}
\usepackage{xspace}
\usepackage{microtype}
\usepackage{enumitem}

\title{KnowAct: Evaluating Functional Theory of Mind through\\Active Knowledge-State Diagnosis}

% Authors must not appear in the submitted version. They should be hidden
% as long as the \iclrfinalcopy macro remains commented out below.
\author{Anonymous Authors}

\newcommand{\knowact}{\textsc{KnowAct}\xspace}
\newcommand{\KG}{\mathcal{G}}
\newcommand{\Vnodes}{\mathcal{V}}
\newcommand{\Eedges}{\mathcal{E}}
\newcommand{\Mtrue}{\mathcal{M}^{\star}}
\newcommand{\Mhat}{\widehat{\mathcal{M}}}
\newcommand{\history}{\mathcal{H}}
\newcommand{\unknown}{\textsc{unknown}}

% \iclrfinalcopy % Uncomment for camera-ready version, but NOT for submission.
\begin{document}

\maketitle

\begin{abstract}
An agent cannot adapt well to a user merely by answering the current request; it must also decide what it still needs to learn about that user. This creates a measurement problem. Static Theory-of-Mind benchmarks test inference from supplied evidence, whereas interactive benchmarks typically score downstream task or social outcomes. The former removes evidence acquisition from the tested policy, and the latter can reward success without an accurate user model.

We argue that functional user modeling should instead be evaluated as active latent-state identification: the target state must be explicit, acquiring evidence must consume a limited budget, and success must be attributable to visible interaction rather than hidden-state leakage or a favorable prior. \knowact instantiates these requirements for knowledge-grounded interaction. A public domain graph defines a shared hypothesis space; a hidden user knowledge map defines the reconstruction target; and a tested agent chooses diagnostic questions before submitting a full-map estimate.

The benchmark separates simulator, agent, and evaluation contexts to address the tension between natural interaction and experimental control. A de-identified answer blueprint constrains what the simulator may express, while deterministic node-level scoring avoids an evaluator-model confound. The initial domain is statistical learning, grounded in an authoritative textbook. This paper develops the benchmark's measurement argument and validation protocol; empirical validation remains future work.
\end{abstract}

\section{Introduction}

In collaborative interaction, the most informative action is often not an answer but a question.

A tutor locates the boundary of a learner's understanding. A research assistant discovers accepted
assumptions. A decision-support agent identifies the uncertainty blocking progress.

In each case, assistance depends on an internal estimate of the user and an action that improves it.

This is more demanding than describing a mental state after its evidence has been supplied.

The agent must decide which uncertainty matters, acquire evidence under a budget, revise its estimate,
and let that estimate change its next action.

User modeling is therefore part of the control loop, not an annotation attached to a completed dialogue.

Existing evaluations illuminate parts of this problem but do not identify the full loop.

Machine Theory-of-Mind (ToM) benchmarks infer beliefs, intentions, or emotions from fixed narratives
and conversations \citep{le2019revisiting,gandhi2023bigtom,kim2023fantom,xu2024opentom,chen2024tombench}.

Their controlled questions make inference measurable, yet the benchmark chooses the evidence. The model
is not accountable for an uninformative question or unresolved uncertainty.

Interactive social-agent environments restore action, but usually score goal completion, persuasion, or
task success \citep{zhou2024sotopia,mou2025agentsense,hwang2026infusing,wang2026tomelp,zhou2026tomswe}.

These outcomes combine user modeling with planning, domain competence, language generation, and
strategy. Success may reflect an accurate user model, a generic policy, or environmental shortcuts.

The inference--application gap in explicit ToM evaluation makes this ambiguity consequential
\citep{gu2024simpletom}.

The resulting gap is one of \emph{measurement}, not simply coverage.

Static tasks remove evidence acquisition. Interactive tasks restore action but often hide the latent
user state from the evaluator.

Neither setting tests whether an agent improves an explicit user model by choosing what evidence to
request.

We study the following question:

\begin{quote}
\emph{Can an agent actively infer a user's hidden knowledge state through limited natural-language interaction, and use that evolving estimate to choose subsequent diagnostic actions?}
\end{quote}

We call the target behavior \emph{functional ToM for knowledge-grounded interaction}.

\emph{Functional} means that a user model earns credit through an infer--update--act loop, not through a
claim that the model understands the user.

\emph{Knowledge-grounded} narrows the latent state to reviewed, diagnosable concepts. The definition
identifies useful behavior without asserting a human-like cognitive mechanism.

A valid test requires three commitments.

First, the latent state must be explicit and directly scoreable. Second, evidence must be selected under
a finite budget. Third, the interaction must be auditable.

Without these commitments, task success, exhaustive questions, leakage, or simulator artifacts can
masquerade as diagnosis.

\knowact turns these commitments into a controlled benchmark.

A public graph defines the shared hypothesis space, while a hidden user map defines the target. The
agent asks diagnostic questions under a fixed turn budget and submits a prediction for every node.

The \sagesim simulator grounds a question before reading hidden state, abstracts only licensed
epistemic content, and realizes it separately. The final map is scored deterministically.

This design makes competing explanations testable.

Fixed and random policies test whether interaction alone is sufficient. Working-map and graph ablations
test whether explicit state guides action.

Simulator variations test proxy dependence. Evidence references separate dialogue-supported
reconstruction from accurate guesses.

It also creates a technical agent problem that the benchmark alone does not solve.

The agent must represent node uncertainty, ground open answers in rubrics, use graph relations without
copying mastery, value the next probe, and decide when evidence is not worth its cost.

We propose MapProbe as an explicit evidence--belief--probe loop. It exposes evidence, beliefs, target
plans, and stopping decisions so a final score can be attributed to a mechanism.

Our contributions are therefore design and protocol contributions rather than empirical claims at this stage:

\begin{itemize}[leftmargin=1.5em]
    \item We formulate active knowledge-state diagnosis as an identifiable target for functional user modeling, with \sagesim as a scoped and independently testable observation model.
    \item We specify MapProbe, a graph-aware, evidence-backed reconstruction agent whose belief update, diagnostic target policy, question realization, and stopping rule are independently inspectable and ablatable.
    \item We predefine a paired evaluation that separates active acquisition, explicit state use, graph use, simulator robustness, repeated reliability, and human transfer instead of collapsing them into one leaderboard score.
\end{itemize}

The visibility boundary, full-map shell, and simple prompted agent baseline are implemented.

An experimental ECDA slice implements node-marginal belief updates and multi-candidate question
utility. The full MapProbe design and comparative experiments remain incomplete or unrun.

The paper therefore does not provide evidence that an agent exhibits the target behavior. This
distinction is deliberate: \knowact makes the claim falsifiable before it is made.

\section{Related Work}
\label{sec:related}

\subsection{From Mental-State Inference to Functional Use}

Machine-ToM benchmarks have made latent-state inference increasingly controlled and difficult. ToMi exposes shortcuts in false-belief question answering \citep{le2019revisiting}; BigToM uses causal templates \citep{gandhi2023bigtom}; and FANToM tests consistency across questions that depend on the same information asymmetry \citep{kim2023fantom}. Hi-ToM, OpenToM, and ToMBench broaden reasoning order, narrative variation, and the taxonomy of social-cognitive abilities \citep{wu2023hitom,xu2024opentom,chen2024tombench}.

These advances strengthen the validity of inference from \emph{supplied} evidence. They do not make evidence acquisition part of the tested policy. SimpleToM identifies a related boundary: explicitly inferring a mental state does not guarantee that a model will use it when predicting or judging behavior \citep{gu2024simpletom}. \knowact takes this boundary as a design requirement. A user model must affect the next diagnostic action, and the result of that process must remain directly scoreable.

Interactive social benchmarks test behavior in richer loops. SOTOPIA and AgentSense evaluate social goals, strategic communication, and private-information reasoning in multi-turn scenarios \citep{zhou2024sotopia,mou2025agentsense}. Applied-ToM agents explicitly insert mental-state hypotheses into social, persuasive, or software-engineering policies \citep{hwang2026infusing,wang2026tomelp,zhou2026tomswe}.

Their question is whether user modeling improves an external outcome. Ours is whether the user model itself is acquired accurately and adaptively. These are complementary claims: downstream success establishes usefulness, whereas direct reconstruction establishes what was learned about the user. Broader agent evaluation likewise shows that single-turn ability and final success can obscure multi-turn and trajectory-level differences \citep{wang2024mint,ma2024agentboard}.

\subsection{Passive User Modeling and Active Diagnosis}

Personalization benchmarks infer preferences and traits from user histories. LaMP evaluates whether profile information improves personalized outputs \citep{salemi2024lamp}; PERSONAMEM asks whether models track profiles that evolve across long interaction histories \citep{jiang2025personamem}. In both cases, the history is evidence to be interpreted. The tested model does not decide which observation would most reduce uncertainty about the user.

Knowledge tracing provides the closest formal analogue. Bayesian and neural tracing methods estimate latent mastery from sequences of item responses \citep{corbett1994knowledgetracing,piech2015dkt,ghosh2020akt}. LongTutor further separates historical evidence acquisition, knowledge diagnosis, and adaptive teaching, revealing that strength at one stage need not transfer to the next \citep{li2026longtutor}.

\knowact changes who controls the evidence. The tested agent asks open-ended questions and is scored on the full latent state after a bounded number of observations. The benchmark is therefore an active diagnosis problem rather than next-response prediction over a fixed curriculum.

\subsection{LLM-Based User Simulation}

User simulation trades experimental control for behavioral fidelity. Agenda-based simulators are reproducible but narrow; LLM simulators are expressive but can invent goals, leak hidden state, or impose model-specific artifacts. DuetSim separates generation and verification to improve response quality \citep{luo2024duetsim}. SimulatorArena instead evaluates the proxy assumption directly, comparing simulated behavior and assistant rankings with human interaction \citep{dou2025simulatorarena}. Earlier simulator research makes the same conceptual point: fluency is not fidelity to task-relevant behavior \citep{shi2019usersimulators}.

\knowact does not treat an LLM simulator as ground truth about human conversation. It uses simulation to hold a reviewed latent state fixed while agents choose evidence. Structural access controls limit what can be expressed, deterministic scoring prevents the simulator from judging success, and matched human episodes test whether conclusions transfer beyond the proxy. Simulator validation is therefore part of the benchmark claim, not an optional naturalness check.

\begin{table}[t]
\caption{Positioning of \knowact relative to adjacent evaluation paradigms. ``Active probes'' means the tested model chooses what evidence to request.}
\label{tab:positioning}
\centering
\small
\begin{tabularx}{\linewidth}{>{\raggedright\arraybackslash}Xcccc}
\toprule
\textbf{Paradigm} & \textbf{Interactive} & \shortstack{\textbf{Explicit latent}\\\textbf{user state}} & \shortstack{\textbf{Active}\\\textbf{probes}} & \shortstack{\textbf{Direct state}\\\textbf{score}} \\
\midrule
Narrative / conversational ToM & Limited & Partial & No & Usually QA \\
Social-agent arenas & Yes & Usually no & Yes & No \\
Personalization from histories & Usually no & Profile / preference & No & Sometimes \\
Knowledge tracing & Sequential & Mastery & Usually no & Prediction-based \\
\textbf{\knowact} & \textbf{Yes} & \textbf{Knowledge map} & \textbf{Yes} & \textbf{Yes} \\
\bottomrule
\end{tabularx}
\end{table}


\section{Problem Formulation}
\label{sec:formulation}

\subsection{Domain Knowledge Graph}

A benchmark domain is represented by an authored knowledge graph
\begin{equation}
    \KG = (\Vnodes, \Eedges),
\end{equation}
where each node $v_i \in \Vnodes$ is a diagnosable knowledge unit and each edge $e_{ij} \in \Eedges$ records a reviewed domain relation. The current relation vocabulary includes \texttt{part\_of}, \texttt{prerequisite\_for}, \texttt{supports}, and \texttt{contrasts\_with}. Each node contains a definition, source locators, a diagnostic goal, mastery-level rubrics, diagnostic signals, and simulator guidance. Edges organize the domain and may guide exploration, but they do not themselves encode user state.

\subsection{Hidden User Knowledge Map}

For a graph with $N = |\Vnodes|$ nodes, a hidden user knowledge map is
\begin{equation}
    \Mtrue = \{s_i^{\star}\}_{i=1}^{N},
\end{equation}
where each state $s_i^{\star}$ includes a discrete mastery level $m_i^{\star} \in \{0,1,\ldots,5\}$, evidence references, misconceptions, and unknowns. The six mastery levels are node-specific and are grounded in authored rubrics rather than interpreted as a universal psychometric scale. A map can additionally be paired with a profile context that affects conversational style, but profile information is not itself part of the primary mastery score.

\subsection{Bounded Diagnostic Episode}

An evaluation episode is specified by
\begin{equation}
    \mathcal{E}p = (\KG, \Mtrue, \mathcal{Z}, T, \rho, \sigma),
\end{equation}
where $\mathcal{Z}$ is simulator-only evidence, $T$ is the maximum number of interaction turns, $\rho$ is the interaction rule, and $\sigma$ is the scoring profile. The initial interaction rule permits one primary diagnostic question per turn. At turn $t$, the tested agent observes the public graph, the visible dialogue history $\history_{t-1}$, and its own working map $\Mhat_{t-1}$. It selects a diagnostic question
\begin{equation}
    q_t \sim \pi_{\theta}(\KG,\history_{t-1},\Mhat_{t-1}),
\end{equation}
receives a simulator answer $a_t$, and updates its estimate
\begin{equation}
    \Mhat_t = U_{\theta}(\Mhat_{t-1}, q_t, a_t).
\end{equation}
After at most $T$ turns, the agent submits a full-graph reconstruction $\Mhat_T$. The hidden map, evidence, simulator traces, and scoring internals are never visible to the tested agent.

\paragraph{Functional ToM in \knowact.}
We define functional ToM operationally as the ability to (i) infer a structured user state from interaction evidence, (ii) maintain and revise that state over time, and (iii) select subsequent actions using the maintained state. The benchmark does not require a particular internal cognitive architecture; it evaluates externally observable behavior under a controlled interface.

\section{The \knowact Benchmark}
\label{sec:benchmark}

\begin{figure}[t]
\centering
\fbox{\begin{minipage}{0.94\linewidth}
\centering
\textbf{Authoritative source} $\rightarrow$ \textbf{Reviewed knowledge graph}\\
$\downarrow$\\[-0.2em]
\textbf{Reviewed hidden knowledge map + evidence}\\
$\downarrow$\\[-0.2em]
\textbf{Simulator} $\leftrightarrow$ \textbf{Tested diagnostic agent}\\
$\downarrow$\\[-0.2em]
\textbf{Final reconstructed map} $\rightarrow$ \textbf{Deterministic scoring}
\end{minipage}}
\caption{High-level \knowact pipeline. Benchmark authoring constructs reviewed hidden artifacts; the tested agent interacts only through the visible episode interface.}
\label{fig:pipeline}
\end{figure}

\subsection{Source-Grounded Graph Authoring}

The benchmark begins with an authoritative domain source. In the initial version, the domain is statistical learning and the selected source is \emph{An Introduction to Statistical Learning with Applications in Python} \citep{james2023islp}. Source material is parsed into ordered segments. A graph-authoring workflow then performs four stages:

\begin{enumerate}[leftmargin=1.7em]
    \item \textbf{Node extraction.} Segment-level model calls propose diagnosable knowledge units and attach source locators and grounding notes.
    \item \textbf{Reconciliation.} Candidate nodes are merged, split, and deduplicated across segments while preserving provenance.
    \item \textbf{Diagnostic enrichment.} Each node receives a definition, diagnostic goal, six-level rubric, diagnostic signals, and simulator behavior guidance.
    \item \textbf{Edge proposal and review.} Candidate relations are proposed using the complete node set and then manually reviewed before promotion.
\end{enumerate}

The authoring model is therefore a proposal mechanism rather than the source of truth. Candidate artifacts are reviewed and promoted into immutable authored graph files. Requiring source locators reduces unconstrained ontology generation from model memory and supports later audit of node definitions and rubrics.

\subsection{Knowledge-Map Authoring}

A hidden map assigns one reviewed state to every node in the authored graph. Profile context, background, preferences, and learning goals can be used to generate a coherent candidate user, but the scored target remains the node-level knowledge map. Evidence records make each state behaviorally realizable: they may specify what the user can explain, where an explanation breaks down, which confusion is likely to appear, or what example the user can and cannot produce.

The map-authoring process separates three artifacts:
\begin{itemize}[leftmargin=1.5em]
    \item the \emph{static state} used as ground truth for scoring;
    \item simulator-only \emph{evidence} that supports natural answers;
    \item optional \emph{profile context} that controls wording style but cannot override the state or add unsupported personal facts.
\end{itemize}
Candidate maps undergo human verification for internal consistency, plausibility, graph coverage, and evaluability before they are eligible for episodes.

\subsection{Immutable Episode Manifests}

An episode manifest binds the benchmark domain, graph version, hidden map identifier, maximum turns, interaction rule, scoring profile, and pinned model/provider configuration. This prevents accidental changes in graph, map, model, or scoring settings across repeated runs. The runtime maintains separate contexts:

\begin{itemize}[leftmargin=1.5em]
    \item \textbf{Tested-agent context:} authored graph, interaction rules, visible transcript, and the agent's own working map.
    \item \textbf{Simulator context:} the current grounded nodes, their hidden states and evidence, visible dialogue needed for continuity, and optional style context.
    \item \textbf{Evaluation context:} hidden map, final reconstruction, and deterministic scoring configuration.
\end{itemize}

This separation is central to the benchmark: it makes the tested agent's information boundary explicit and reduces hidden-state leakage through shared prompts or debug traces.

\section{User Simulator}
\label{sec:simulator}

The simulator must satisfy two competing requirements: answers should be natural enough to support open-ended diagnosis, yet constrained enough to remain faithful to the hidden map. \knowact implements the simulator as a staged information boundary rather than a single unconstrained role-playing prompt.

\subsection{Question Grounding}

Given the current question, visible graph, and limited visible dialogue, a grounder identifies the directly targeted node set $S_t \subseteq \Vnodes$. It also detects whether the turn contains multiple independent questions, requests hidden benchmark labels, or cannot be grounded. Crucially, grounding has no access to the hidden map or hidden evidence. This prevents the simulator from interpreting a vague question in whichever way would most conveniently reveal hidden state.

A question may target multiple nodes only when it forms one integrated diagnostic probe, such as asking for a comparison between two methods. Multiple independent questions are rejected or clarified because they would violate the one-question-per-turn budget.

\subsection{Hidden Context Selection and Answer Policy}

Only after grounding does the context builder retrieve hidden state and evidence, and only for nodes in $S_t$. Neighboring nodes are not exposed merely because they are connected in the graph. An answer policy then maps the selected hidden information to a structured \emph{simulator answer blueprint}. The blueprint records a response mode, primary stance, answer shape, and de-identified content units such as:

\begin{itemize}[leftmargin=1.5em]
    \item a claim the user can express;
    \item a boundary of partial or fragile understanding;
    \item a misconception or uncertainty that should surface;
    \item evidence-supported cues that may be used;
    \item claims that would overstate the user's ability.
\end{itemize}

The blueprint intentionally excludes mastery labels, node identifiers, evidence identifiers, map identifiers, user identifiers, and scoring fields. It is the sole content interface to natural-language generation.

\subsection{Answer Generation and Leakage Control}

The answer generator receives the blueprint, visible dialogue required for continuity, and an optional style hint. It never receives the raw hidden map. Generation is instructed to speak in the first person, avoid benchmark schema language, and preserve uncertainty or misconceptions rather than converting all states into binary know/do-not-know answers. Failed or malformed generations use bounded retries and safe non-leaking fallbacks.

This design does not guarantee that an LLM simulator perfectly matches human behavior. It instead provides a traceable contract: grounding determines which state may be read, policy determines what content may be expressed, and generation controls only surface realization. Hidden policy traces remain available to benchmark authors for audit but are not part of the visible transcript.

\section{Tested-Agent Protocol}
\label{sec:agent}

\subsection{Full-Graph Working Map}

The tested agent maintains a working assessment for every graph node. Each node record contains:
\begin{equation}
    \hat{s}_{i,t} = (\hat{m}_{i,t}, c_{i,t}, n_{i,t}, R_{i,t}),
\end{equation}
where $\hat{m}_{i,t} \in \{\unknown,0,\ldots,5\}$ is assessed mastery, $c_{i,t}$ is diagnostic confidence, $n_{i,t}$ is an optional assessment note, and $R_{i,t}$ is the set of visible turn identifiers supporting the assessment. All nodes begin as \unknown. A non-unknown assessment should cite at least one visible observation; unsupported estimates can still be submitted but are reported separately.

A full-graph map serves two purposes. First, it forces the agent to represent uncertainty explicitly instead of silently omitting unvisited concepts. Second, it supports action selection over the entire domain: the agent can compare assessed mastery, confidence, graph position, and evidence coverage when deciding where to probe next.

\subsection{Semantic Action Interface}

Rather than granting arbitrary state mutation, the runtime exposes a small set of semantic actions:

\begin{itemize}[leftmargin=1.5em]
    \item read the authored episode graph;
    \item read the current working map;
    \item update one node assessment with visible support;
    \item ask one diagnostic question;
    \item finalize the reconstructed map.
\end{itemize}

The interface makes benchmark-relevant behavior observable. A run trace records which nodes were updated, which observations supported those updates, how questions changed over time, and when the agent finalized. It also prevents the evaluated policy from bypassing the intended loop through generic access to hidden artifacts or unrestricted storage.

\subsection{Diagnostic Loop}

A generic tested agent follows four recurring steps:

\begin{enumerate}[leftmargin=1.7em]
    \item \textbf{Interpret evidence:} map the latest answer to one or more visible graph concepts.
    \item \textbf{Update beliefs:} revise mastery, confidence, notes, and supporting observations.
    \item \textbf{Select a probe:} choose a node or relation whose answer is expected to reduce decision-relevant uncertainty.
    \item \textbf{Ask or finalize:} generate one grounded diagnostic question or submit the final map when the budget is exhausted.
\end{enumerate}

\knowact does not prescribe how uncertainty or information value must be computed. A baseline may use a single LLM prompt, while more structured agents may use entropy estimates, graph-aware coverage, Bayesian updates, planning, or learned policies. The shared working-map and action protocol makes these designs comparable.

\section{Evaluation}
\label{sec:evaluation}

\subsection{Primary Reconstruction Metric}

The primary metric is episode mastery distance. Let the final predicted mastery for node $i$ be $\hat{m}_i$. For submitted levels $0$ through $5$, the node loss is squared ordinal distance:
\begin{equation}
    d(\hat{m}_i,m_i^{\star}) = (\hat{m}_i-m_i^{\star})^2.
\end{equation}
An \unknown prediction receives the maximum fixed penalty of $36$, corresponding to a distance larger than any error between two valid levels. The episode score is
\begin{equation}
    D_{\mathrm{mastery}} = \frac{1}{N}\sum_{i=1}^{N} d(\hat{m}_i,m_i^{\star}),
\end{equation}
where lower is better. The score is deterministic and requires no evaluator model.

The ordinal loss rewards near misses and penalizes severe over- or under-estimation more strongly than exact-match accuracy alone. It should nevertheless be interpreted relative to each node's authored rubric; the numeric levels are a benchmark encoding, not an assumption that all conceptual differences are psychometrically equal.

\subsection{Supporting Metrics}

The benchmark can report the following secondary metrics without changing the primary objective:
\begin{itemize}[leftmargin=1.5em]
    \item exact mastery match rate and per-node signed error;
    \item missing prediction rate;
    \item unsupported inference rate, based on absent visible evidence references;
    \item graph and rubric coverage under a fixed turn budget;
    \item reconstruction quality as a function of turn count;
    \item redundant or ungrounded question rate;
    \item information gain per turn for agents that expose calibrated beliefs.
\end{itemize}
Misconception and free-text unknown fields are retained for qualitative analysis but are not included in the initial automatic score, avoiding an additional subjective judge.

\subsection{Baseline Families}

The initial benchmark supports three baseline families:
\begin{itemize}[leftmargin=1.5em]
    \item \textbf{Fixed-question agent:} follows a predefined diagnostic order and never adapts question choice to answers.
    \item \textbf{Random-question agent:} samples valid diagnostic targets or questions under episode constraints.
    \item \textbf{Simple LLM agent:} observes the graph, working map, and visible dialogue, then directly chooses the next question and updates the map through the semantic interface.
\end{itemize}
These baselines separate the value of interaction from the value of adaptive user modeling. Future methods can add explicit uncertainty estimation, graph-aware planning, information-gain objectives, specialized user-state modules, or multi-agent decomposition while using the same simulator and scoring protocol.

\section{Experiments}
\label{sec:experiments}
% Intentionally left blank in Draft v1.
% Planned subsections:
% - Research questions
% - Benchmark instances and human review statistics
% - Models and agent baselines
% - Main reconstruction results
% - Turn-efficiency analysis
% - Ablations of working-map structure, graph information, and simulator boundaries
% - Simulator fidelity study with human participants
% - Error analysis

\section{Discussion}
\label{sec:discussion}

\subsection{What the Benchmark Measures}

A low reconstruction error alone is not sufficient evidence of functional ToM if it is achieved through leakage or a fixed prior. \knowact therefore couples endpoint accuracy with process observability. The tested agent must operate under a strict visibility boundary, questions consume a finite budget, map updates cite visible turns, and fixed or random policies provide non-adaptive controls. Together, these constraints aim to distinguish three abilities that are often conflated:

\begin{enumerate}[leftmargin=1.7em]
    \item \emph{state inference}: interpreting what an answer implies about the user;
    \item \emph{state maintenance}: integrating evidence into a coherent map over time;
    \item \emph{state-conditioned action}: using the current map to select the next probe.
\end{enumerate}

The benchmark is deliberately diagnostic rather than task-complete. It does not yet test whether the recovered map improves teaching, recommendation, or collaborative problem solving. Those downstream tasks can be layered on top once the state-reconstruction loop is validated.

\subsection{Relation to Active Learning and Assessment}

The interaction can be viewed as active information acquisition over a structured latent state. Unlike pool-based active learning, however, the action is a natural-language question whose interpretation and answer are mediated by a conversational simulator. Unlike computerized adaptive testing, the item set is not necessarily fixed and calibrated in advance; agents may generate integrated conceptual probes. This flexibility motivates the benchmark but also increases variance, making grounding, evidence audit, and simulator validation important.

\subsection{Limitations}

\paragraph{Simulator fidelity.}
A simulator conditioned on a reviewed map can be internally consistent without matching how real users express partial knowledge. Human studies are required to measure response similarity, diagnostic usefulness, and whether agent rankings transfer from simulated to human interaction.

\paragraph{Construction circularity.}
If related model families author the graph, generate the hidden map, simulate the user, and act as the tested agent, correlated biases may make the benchmark easier or distort rankings. Source grounding, human promotion, model-family diversity, rule-based fixtures, and human validation can reduce but not eliminate this risk.

\paragraph{Simplified knowledge state.}
A static L0--L5 state per node cannot capture all aspects of human knowledge, including context-dependent performance, confidence, procedural fluency, forgetting, and learning during interaction. The initial benchmark freezes the hidden state during an episode to make reconstruction identifiable. Dynamic learning is an important extension.

\paragraph{Single-domain initial scope.}
The first release focuses on statistical learning. Cross-domain evaluation is needed to determine whether agent policies transfer beyond one graph topology, source style, and concept distribution.

\paragraph{Operational rather than cognitive claims.}
Success on \knowact demonstrates behavior compatible with structured user modeling under benchmark constraints. It does not establish that a model possesses human-like beliefs, consciousness, or a psychologically faithful ToM mechanism.

\section{Ethical Considerations}

User modeling can improve accessibility, tutoring, and collaborative assistance, but it can also enable manipulation, surveillance, or unjustified inference about individuals. \knowact uses synthetic or explicitly reviewed benchmark profiles and should not be interpreted as permission to infer sensitive attributes from real users without consent. Systems derived from this work should minimize retained personal data, expose uncertainty, allow users to inspect or correct inferred states, and avoid using knowledge-state estimates for high-stakes decisions without appropriate human oversight. Benchmark releases should also document the provenance and licenses of source material and avoid embedding identifiable information in profile contexts or evidence records.

\section{Conclusion}

We presented \knowact, a benchmark framework for evaluating functional Theory of Mind through active knowledge-state diagnosis. The framework gives a tested agent a reviewed domain graph but hides a user-specific knowledge map. Under a bounded interaction budget, the agent must ask diagnostic questions, update an evidence-backed working map, and reconstruct the hidden state. A visibility-separated simulator and deterministic scoring protocol make the process auditable and reproducible. By isolating active user-state acquisition from broader social or task success, \knowact provides a controlled foundation for comparing agent architectures that infer, maintain, and act on models of users. The next step is empirical validation across agent baselines, simulator designs, human participants, and additional knowledge domains.

\bibliography{references}
\bibliographystyle{iclr2026_conference}

\appendix

\section{Implementation-Oriented Artifact Schemas}
\label{app:schemas}

The following abbreviated schemas summarize the current benchmark contracts. Exact field names may evolve during implementation, but the separation between public graph, hidden map, visible observations, and final reconstruction should remain stable.

\subsection{Knowledge Node}
\begin{verbatim}
{
  "id": "active_learning",
  "name": "Active Learning",
  "type": "concept",
  "definition": "...",
  "source_locators": [...],
  "diagnostic_goal": "...",
  "levels": {"L0": "...", ..., "L5": "..."},
  "diagnostic_signals": [...],
  "simulator_behavior": "..."
}
\end{verbatim}

\subsection{Hidden User State}
\begin{verbatim}
{
  "node_id": "active_learning",
  "mastery_level": "L2",
  "evidence_refs": ["ev_104"],
  "misconceptions": [...],
  "unknowns": [...]
}
\end{verbatim}

\subsection{Agent Working Assessment}
\begin{verbatim}
{
  "node_id": "active_learning",
  "assessed_mastery_level": "L2",
  "diagnostic_confidence": "medium",
  "assessment_note": "Understands the motivation but not
                      query-strategy trade-offs.",
  "supporting_turn_ids": ["turn_03"]
}
\end{verbatim}

\section{Planned Empirical Questions}
\label{app:planned}

The empirical study should be designed around questions that directly test the benchmark claim rather than only comparing foundation-model capability:

\begin{enumerate}[leftmargin=1.7em]
    \item Do adaptive agents reconstruct hidden maps more accurately than fixed and random diagnostic policies under equal turn budgets?
    \item Does an explicit working map improve reconstruction and question efficiency relative to transcript-only prompting?
    \item Do graph relations improve target selection, or do they induce unsupported propagation between neighboring nodes?
    \item How sensitive are rankings to simulator model, answer-policy implementation, question grounding, and profile style?
    \item How well do simulator answers and agent rankings agree with matched human-user episodes?
\end{enumerate}

\end{document}
